% Generated draft: requires mathematical and source review.
\documentclass[11pt]{article}
\usepackage[margin=25mm]{geometry}
\usepackage{fontspec}
\setmainfont{Latin Modern Roman}
\usepackage{amsmath,amssymb,mathtools,unicode-math}
\setmathfont{Latin Modern Math}
\usepackage{enumitem,needspace,xurl,hyperref}
\setlength{\parindent}{0pt}
\setlength{\parskip}{5pt}
\setlength{\emergencystretch}{3em}
\newcommand{\N}{\mathbb{N}}
\newcommand{\Z}{\mathbb{Z}}
\newcommand{\Q}{\mathbb{Q}}
\newcommand{\R}{\mathbb{R}}
\newcommand{\C}{\mathbb{C}}
\newcommand{\F}{\mathbb{F}}
\newcommand{\A}{\mathbb{A}}
\newcommand{\PP}{\mathbb{P}}
\newcommand{\E}{\mathbb{E}}
\newcommand{\eps}{\varepsilon}
\newcommand{\dd}{\,\mathrm{d}}
\begin{document}
\section*{UnsolvedMath ID: 2303013}
\textbf{Source problem code:} catalogue entry
\subsection{Definitions and mathematical statement}
For a subharmonic function $u$ on $\mathbb R^m$, its mass means its Riesz measure, a fixed positive normalization of the distributional Laplacian $\Delta u$. Saying that all mass lies on $S$ means $\operatorname{supp}(\Delta u)\subseteq S$; it does not say that $u$ itself vanishes off $S$. Let $N(r,u)$ and $T(r)$ be the higher-dimensional Nevanlinna counting and characteristic functions intended by the source, and put
\[
\delta(u)=1-\limsup_{r\to\infty}\frac{N(r,u)}{T(r)},\qquad
\rho(u)=\limsup_{r\to\infty}\frac{\log^+T(r)}{\log r}
\]
when the characteristic is nontrivial. The source refers to analogues of planar counting functions but gives no formulas or normalizations; these must be fixed for numerical sharp bounds.

\textbf{Questions (D. Shea).} For order $\rho>1$, determine the corresponding sharp upper bound on $\delta(u)$, analogous to the supplied result for $\rho<1$. In addition, if the Riesz mass is supported on a ray, can one find a uniform lower bound $A(\rho)$ on $\delta(u)$? More generally, for a fixed suitable lower-dimensional set $S\subset\mathbb R^m$, find a threshold $\rho_0(S)$ and a bound, written with suppressed dependencies as $A(\rho)$ in the source, such that
\[
\delta(u)\geq A_{m,S}(\rho)
\]
for every eligible subharmonic $u$ of order $\rho>\rho_0(S)$ with Riesz measure supported on $S$. The conjectured high-order behavior is
\[
\lim_{\rho\to\infty}A_{m,S}(\rho)=1
\]
for fixed $m$ and $S$, in the range where the class is nonempty. This asks for a bound uniform over the functions of each order, not merely a statement about one fixed function.

The source reports a sharp upper bound for $\rho<1$, attained with mass on a ray, and attributes the planar high-order lower-bound result to Hellerstein and Shea. Its research summary describes the higher-dimensional $m>2$, $\rho>1$ problem as open as of 2018, with no recent resolution located.
\subsection{Short English statement}
Study deficiency bounds for high-order subharmonic functions, especially when their Riesz mass lies on a ray or another fixed lower-dimensional set. The proposed lower bound should approach one as the order tends to infinity.
\subsection{Sources}
\begin{itemize}
\item Hayman, \emph{Research Problems in Function Theory} (2018), Problem 3.13, attributed to D. Shea, \url{https://arxiv.org/abs/1809.07200}; source-TeX extraction, accessed 2026-07-29. The 2018 update was retained for status triage.
\item Hayman, higher-dimensional Nevanlinna definitions; Hayman and Kennedy, the low-order sharp bound; Hellerstein and Shea, the planar high-order result. The extracted references provide no titles, dates, or theorem locations.
\item Local registry \texttt{data/unsolvedmath/problems.json}, permanent ID $2303013$, source code catalogue entry, \emph{Research Problems in Function Theory --- Problem 3.13}. The record labels this open and OPEN-TRIAGE, with NEEDS\_REVIEW source-status and rights qualifications.
\end{itemize}
\end{document}
